% Ad Familiares 11.6 (ad-familiares-11-06) --- English translation
% Translated by Claude (Anthropic) from the Latin (Perseus).
% Style guide: STYLE.md.

\ciceroLetterOpener{M. CICERO S. D. D. BRVTO IMP. COS. DESIG.}

\ciceroSection{1} Our friend Lupus, when he had come to Rome from Mutina on the sixth day, met me the next day in the morning; he set out your charges to me with the utmost diligence, and handed over your letter. In that you commend your standing to me, I think that at the same time you are commending my standing to me --- which, by Hercules, I do not hold dearer than your own. For which reason you will do me a very welcome service if you hold it for certain that neither my counsel nor my zeal will fail your praise at any point.

\ciceroSection{2} When the tribunes of the plebs had given notice that the Senate should meet on the twentieth of December, and had it in mind to bring the matter of a guard for the consuls-elect, although I had determined not to come to the Senate before the Kalends of January, nonetheless --- since on that very day your edict had been put up --- I judged it impious either that the Senate should sit in such a way that there should be silence about your divine services to the state (which would have happened, had I not come), or even, if anything honourable should be said about you, that I should not be present. And so I came to the Senate in the morning.

\ciceroSection{3} Once this had been noticed, the senators assembled in great numbers. What I did about you in the Senate, what I said in the great public assembly, I prefer that you should learn from others' letters; let me have you persuaded only of this --- that I shall always undertake and defend with the utmost zeal everything that shall bear on increasing your standing, which in itself is at the height. And although I see that in this I shall be acting along with many, nonetheless I shall make a bid for the first place in this business.
